\documentclass[10pt,a4paper]{amsart}
\usepackage[margin=19mm]{geometry}
\usepackage{amsmath,amssymb,mathtools}
\usepackage[scr=rsfs,cal=euler]{mathalfa}
\usepackage{xfrac}
\usepackage[unicode,hidelinks,bookmarks=false]{hyperref}
\newcommand{\ie}{i.e.\ }
\newcommand{\cf}{cf.\ }
\newcommand{\resp}{respectively\ }
\newcommand{\Subsection}{\subsection}
\providecommand{\nobreakdash}{\nobreak\hbox{-}}
\setlength{\emergencystretch}{2em}
\multlinegap0pt
\usepackage{bm}
\usepackage{bbm}
\usepackage{dsfont}
\usepackage{xspace}
\usepackage{cancel}
\usepackage{mathtools}
\usepackage{amsthm}
\makeatletter
\@ifundefined{theorem}{\newtheorem{theorem}{Theorem}}{}
\@ifundefined{assumption}{\newtheorem{assumption}[theorem]{Assumption}}{}
\@ifundefined{proposition}{\newtheorem{proposition}[theorem]{Proposition}}{}
\makeatother
\newcommand{\dto}{\xrightarrow{d}}
\title[Ball-Codifference Screening for Heavy-Tailed Predictors]{Ball-Codifference Screening for Heavy-Tailed Predictors}
\author{Mohsen Rezapour, Vahed Maroufy}
\date{Source: arXiv:2607.20821v1 (CC BY 4.0)}
\begin{document}
\maketitle

\noindent\emph{Source and licence.} Ball-Codifference Screening for Heavy-Tailed Predictors. Version arXiv:2607.20821v1 (CC BY 4.0), distributed under the Creative Commons Attribution 4.0 International licence, as recorded in the arXiv licence metadata for this exact version and shown on the abstract page https://arxiv.org/abs/2607.20821v1. Full rights notice: licenses/arxiv-2607.20821-CC-BY-4.0.txt.\par\smallskip

\noindent\textit{Editorial scope.} Counted unit: the Proposition whose source label is \texttt{prop:vg\_derivative} in the article (source0.tex lines 686--700, source label \texttt{prop:vg\_derivative}), delivered together with the article's own complete original proof of it (source0.tex lines 702--732, one \texttt{proof} environment of 31 lines). No mathematics is written, repaired or re-derived: statement and proof are the article's own text line for line. Required context retained uncounted: eq:Vg (lines 671--673); ass:vg (lines 682--684). Every definition, display and label of those lines is retained unchanged. Omitted from this package and not redistributed: 1--670, 674--681, 685--685, 701--701, 733--825. Nothing inside a retained excerpt is omitted or altered: every retained body is delivered exactly as the article's lines are written, so no omission receipt of any kind accompanies this package. Citations: the retained excerpts carry no citation command at all, so no bibliography entry of the article is transcribed and no reference is redistributed. Reference adaptation: none was needed and none was made; every reference inside the delivered unit resolves inside it, so no unresolved reference or citation is produced. Printed slips kept literal: no printed word, symbol, hypothesis or number of the counted statement or of its proof was altered. Layout and class: the article's own class is \texttt{article} with packages including geometry, amsmath, amssymb, amsfonts, amsthm, mathtools, booktabs, array, multirow, graphicx, natbib, hyperref, enumitem, placeins; the wrapper is the lane's pinned \texttt{amsart} editorial wrapper at 10pt on A4, which restates the theorem-like environments the retained text needs and adds the article's own macro definitions that the retained text needs (\texttt{\textbackslash dto}), with extra wrapper packages (bm, bbm, dsfont, xspace, cancel, mathtools). The delivered numbering is the unit's own. Assets: no retained span contains a figure environment, a graphics-inclusion command, a PDF-page inclusion, a table or a TikZ picture; no image file is copied into this package, the package declares no asset dependency and no image file is redistributed with it. Licence: version arXiv:2607.20821v1 of this article is distributed under the Creative Commons Attribution 4.0 International licence, as recorded in the arXiv licence metadata for this exact version and shown on the abstract page https://arxiv.org/abs/2607.20821v1. A full rights notice is delivered at licenses/arxiv-2607.20821-CC-BY-4.0.txt. Characters: every retained excerpt is pure ASCII, so the delivered engine, XeLaTeX with its default Latin Modern text font, reports no missing character.
\begin{equation}\label{eq:Vg}
V_g(F)=\int\cdots\int g(x_1,\ldots,x_m)F(dx_1)\cdots F(dx_m),
\end{equation}

\begin{assumption}\label{ass:vg}
The following hold: (i) the integrals in \eqref{eq:Vg} exist for \(F\) and for the empirical distribution \(F_n\); (ii) \(|g|\le M\) for some finite \(M\); (iii) each \(g_{i,F}\) is locally of bounded variation and \(\int |dg_{i,F}|<\infty\); (iv) the integration-by-parts boundary terms vanish; and (v) \(\sqrt n(F_n-F)\dto B^\circ\) in the tangent space.
\end{assumption}

\begin{proposition}\label{prop:vg_derivative}
Under Assumption~\ref{ass:vg}, the functional \(V_g\) is quasi-Hadamard differentiable at \(F\) tangentially to \(\mathbb C^q\), with derivative
\begin{equation}\label{eq:vg_derivative}
\dot V_{g,F}(v)=\sum_{i=1}^m(-1)^q\int v(x-)\,dg_{i,F}(x).
\end{equation}
Consequently,
\[
\sqrt n\{V_g(F_n)-V_g(F)\}\dto \dot V_{g,F}(B^\circ).
\]
If \(B^\circ\) is a centered Gaussian process with covariance function \(\Gamma\), then \(\dot V_{g,F}(B^\circ)\) is normal with variance
\[
\sum_{i,j=1}^m\int\int \Gamma(x,y)\,dg_{i,F}(x)\,dg_{j,F}(y),
\]
whenever the displayed integrals exist.
\end{proposition}

\begin{proof}
Let \(h_n\to0\), let \(v_n\to v\) in the tangent norm, and put \(F_n^*=F+h_nv_n\).  We must show that
\[
\frac{V_g(F_n^*)-V_g(F)}{h_n}\to \dot V_{g,F}(v).
\]
Expanding the product measure in \eqref{eq:Vg} gives
\begin{align}\label{eq:product_expansion}
\frac{V_g(F_n^*)-V_g(F)}{h_n}
&=\sum_{i=1}^m
\int\cdots\int g(x_1,\ldots,x_m)
\prod_{\ell<i}F(dx_\ell)\,v_n(dx_i)\prod_{\ell>i}F(dx_\ell) \\
&\quad +\sum_{r=2}^m h_n^{r-1} Q_{n,r},\nonumber
\end{align}
where \(Q_{n,r}\) is a finite sum of integrals in which \(r\) coordinates are integrated with respect to \(v_n\) and the remaining coordinates with respect to \(F\).  Because \(|g|\le M\) and the total variations of \(v_n\) are bounded along convergent tangent sequences, there is a constant \(C\), independent of \(n\), such that \(|Q_{n,r}|\le C\) for every \(r=2,\ldots,m\).  Hence the second term in \eqref{eq:product_expansion} is \(o(1)\).

For the first-order terms, Fubini's theorem gives
\[
\int\cdots\int g(x_1,\ldots,x_m)
\prod_{\ell<i}F(dx_\ell)\,v_n(dx_i)\prod_{\ell>i}F(dx_\ell)
=\int g_{i,F}(x)\,v_n(dx).
\]
We next show that \(\int g_{i,F}\,dv_n\to \int g_{i,F}\,dv\).  Since \(g_{i,F}\) has finite total variation on compact sets and \(v_n\to v\) uniformly at continuity points, the convergence follows first on rectangles by a standard partition argument.  The tails are controlled by \(\int|dg_{i,F}|<\infty\) and the bounded total variations of \(v_n\).  Therefore
\[
\int g_{i,F}(x)\,v_n(dx)\to \int g_{i,F}(x)\,v(dx).
\]
Finally, repeated one-dimensional integration by parts in the \(q\) coordinates, with the boundary terms equal to zero by Assumption~\ref{ass:vg}, yields
\[
\int g_{i,F}(x)\,v(dx)=(-1)^q\int v(x-)\,dg_{i,F}(x).
\]
Summing over \(i\) proves the derivative formula \eqref{eq:vg_derivative}.  The weak convergence statement follows from the modified functional delta method applied to \(\sqrt n(F_n-F)\dto B^\circ\).  The variance formula follows by linearity of \(\dot V_{g,F}\) and the covariance structure of the Gaussian process.
\end{proof}

\end{document}
